% 原书第203—208页；10.3.3。
\subsection{Granger因果性}
为简化起见，我们先介绍Granger因果性的二元版本。

\paragraph*{二元Granger因果性}
定理10.4说明，在排除即时效应并满足一些温和技术条件时，通过检验式(10.5)，再交换$X$与$Y$的角色检验相应陈述，就能推断汇总图中的箭头是否存在。由此可以区分四种可能的汇总图：$X\quad Y$、$X\to Y$、$X\leftarrow Y$以及$X\rightleftarrows Y$。如果在根据$Y$自身过去预测$Y$时，加入$X$的过去能够改善预测，就推断$X$影响$Y$。形式上，写作
\begin{equation}
 \text{$X$在Granger意义下导致$Y$}
 \quad:\Longleftrightarrow\quad
 Y_t\nindep X_{\mathrm{past}(t)}\mid Y_{\mathrm{past}(t)}.
 \tag{10.6}
\end{equation}
这一思想可以追溯到Wiener［1956，第189—190页］：他认为，如果在利用$Y$自身过去进行预测时，额外考虑$X$能够改善预测，那么$X$就对$Y$有因果影响。图10.6展示了满足定理10.4的典型情形。
\begin{figure}[htbp]
\centering
\begin{tikzpicture}[causal,x=1.45cm,y=1.6cm]
\node[ch10node] (ch10f6-X-0) at (0,1) {$X_{t-6}$};
\node[ch10node] (ch10f6-X-1) at (1,1) {$X_{t-5}$};
\node[ch10node] (ch10f6-X-2) at (2,1) {$X_{t-4}$};
\node[ch10node] (ch10f6-X-3) at (3,1) {$X_{t-3}$};
\node[ch10node] (ch10f6-X-4) at (4,1) {$X_{t-2}$};
\node[ch10node] (ch10f6-X-5) at (5,1) {$X_{t-1}$};
\node[ch10node] (ch10f6-X-6) at (6,1) {$X_{t}$};
\draw[ed] (ch10f6-X-0)--(ch10f6-X-1);
\draw[ed] (ch10f6-X-1)--(ch10f6-X-2);
\draw[ed] (ch10f6-X-2)--(ch10f6-X-3);
\draw[ed] (ch10f6-X-3)--(ch10f6-X-4);
\draw[ed] (ch10f6-X-4)--(ch10f6-X-5);
\draw[ed] (ch10f6-X-5)--(ch10f6-X-6);
\node[ch10node] (ch10f6-Y-0) at (0,0) {$Y_{t-6}$};
\node[ch10node] (ch10f6-Y-1) at (1,0) {$Y_{t-5}$};
\node[ch10node] (ch10f6-Y-2) at (2,0) {$Y_{t-4}$};
\node[ch10node] (ch10f6-Y-3) at (3,0) {$Y_{t-3}$};
\node[ch10node] (ch10f6-Y-4) at (4,0) {$Y_{t-2}$};
\node[ch10node] (ch10f6-Y-5) at (5,0) {$Y_{t-1}$};
\node[ch10node] (ch10f6-Y-6) at (6,0) {$Y_{t}$};
\draw[ed] (ch10f6-Y-0)--(ch10f6-Y-1);
\draw[ed] (ch10f6-Y-1)--(ch10f6-Y-2);
\draw[ed] (ch10f6-Y-2)--(ch10f6-Y-3);
\draw[ed] (ch10f6-Y-3)--(ch10f6-Y-4);
\draw[ed] (ch10f6-Y-4)--(ch10f6-Y-5);
\draw[ed] (ch10f6-Y-5)--(ch10f6-Y-6);
\draw[ed] (ch10f6-X-0)--(ch10f6-Y-1);
\draw[ed] (ch10f6-X-1)--(ch10f6-Y-2);
\draw[ed] (ch10f6-X-2)--(ch10f6-Y-3);
\draw[ed] (ch10f6-X-3)--(ch10f6-Y-4);
\draw[ed] (ch10f6-X-4)--(ch10f6-Y-5);
\draw[ed] (ch10f6-X-5)--(ch10f6-Y-6);
\draw[ed] (ch10f6-X-0)--(ch10f6-Y-2);
\draw[ed] (ch10f6-X-1)--(ch10f6-Y-3);
\draw[ed] (ch10f6-X-2)--(ch10f6-Y-4);
\draw[ed] (ch10f6-X-3)--(ch10f6-Y-5);
\draw[ed] (ch10f6-X-4)--(ch10f6-Y-6);
\draw[ed] (ch10f6-Y-0)--(ch10f6-X-2);
\draw[ed] (ch10f6-Y-1)--(ch10f6-X-3);
\draw[ed] (ch10f6-Y-2)--(ch10f6-X-4);
\draw[ed] (ch10f6-Y-3)--(ch10f6-X-5);
\draw[ed] (ch10f6-Y-4)--(ch10f6-X-6);
\end{tikzpicture}
\caption{Granger因果性能够奏效的典型情形：如果删除所有从$X$到$Y$的箭头，给定$Y$自身过去后，$Y_t$将与$X$的过去条件独立。这里，给定自身过去后，$Y_t$仍依赖$X$的过去。因此，条件(10.4)证明了从$X$到$Y$的影响存在。}
\end{figure}


Granger因果性通常指线性预测。这时比较下面两个线性回归模型：
\begin{equation}
 Y_t=\sum_{i=1}^q a_iY_{t-i}+N_t,
 \tag{10.7}
\end{equation}
\begin{equation}
 Y_t=\sum_{i=1}^q a_iY_{t-i}+\sum_{i=1}^q b_iX_{t-i}+\widetilde N_t,
 \tag{10.8}
\end{equation}
其中假设$(N_t)_{t\in\mathbb Z}$与$(\widetilde N_t)_{t\in\mathbb Z}$各自都是独立同分布时间序列。如果纳入$X$进行预测时，噪声项$\widetilde N_t$的方差显著小于不纳入$X$时的噪声项$N_t$的方差，就推断$X$在Granger意义下导致$Y$。这相当于说，给定$Y_{\mathrm{past}(t)}$，$Y_t$与$X_{\mathrm{past}(t)}$存在非零偏相关。对多元高斯分布，这等价于相依陈述(10.4)。采用非线性回归的各种改进也已得到广泛研究［例如Ancona et al., 2004；Marinazzo et al., 2008］。关于式(10.5)的非参数检验，可参见Diks and Panchenko［2006］及其参考文献。

\textbf{转移熵（transfer entropy）}［Schreiber, 2000］，是在给定$Y$的过去后，度量$Y_t$与$X$的过去之间相依程度的信息论量：
\begin{equation}
 \mathrm{TE}(X\to Y):=I\bigl(Y_t:X_{\mathrm{past}(t)}\mid Y_{\mathrm{past}(t)}\bigr),
 \tag{10.9}
\end{equation}
其中，对任意三个变量集$\mathbf A,\mathbf B,\mathbf C$，$I(\mathbf A:\mathbf B\mid\mathbf C)$表示条件互信息［Cover and Thomas, 1991］，另见附录A。因此，可以估计转移熵，并在其显著大于$0$时推断$X$导致$Y$；这可以看作Granger因果性的一种信息论实现，能够考虑任意非线性影响。由此，很容易想将转移熵用作$X$对$Y$影响\emph{强度}的度量，但下面关于Granger因果性局限性的讨论将说明，这并不合适。

\paragraph*{多元Granger因果性}
像定理10.4那样，假设二元时间序列满足因果充分性，往往并不恰当。Granger［1980］已经讨论过这一问题。因此，当
\[
 X_t^k\nindep X_{\mathrm{past}(t)}^j\mid\mathbf X_{\mathrm{past}(t)}^{-j}
\]
成立时，我们称$X^j$在Granger意义下导致$X^k$。Granger早已强调，恰当地使用Granger因果性，实际上需要以世界上一切相关变量为条件。不过，人们仍经常使用其二元版本，或将它用于一些明显重要的变量未被观测到的情形。若对所得结果作因果解释，这种用法可能产生误导。

\paragraph*{Granger因果性的局限}
与前几章的独立同分布情形一样，违反因果充分性，是因果时间序列分析中的一个严重问题。为说明Granger因果性为何会在因果不充分的多元时间序列中误导我们，先只考虑观测到二元时间序列$(X_t,Y_t)_{t\in\mathbb Z}$的情况。假设$X_t$与$Y_t$都受到潜在时间序列$(Z_t)_{t\in\mathbb Z}$的较早实例的影响。图10.7(a)展示了这一情形：$Z$以一期时滞影响$X$，以两期时滞影响$Y$。假设忠实性成立，则d-分离判据告诉我们
\[
 Y_t\nindep X_{\mathrm{past}(t)}\mid Y_{\mathrm{past}(t)},
\]
而另一方面有
\[
 X_t\indep Y_{\mathrm{past}(t)}\mid X_{\mathrm{past}(t)}.
\]
因此，直接应用Granger因果性会推断$X$导致$Y$，而$Y$不导致$X$。例如，黄油价格与奶酪价格之间就观察到过这种现象。两者都受到牛奶价格的强烈影响，但生产奶酪比生产黄油耗时长得多，因此牛奶价格与奶酪价格之间的时滞更大［Peters et al., 2013，实验10］。不过，Granger因果性的这类失效，正是由于未观测所有相关变量，而Granger本人已明确要求满足这一条件。
\begin{figure}[htbp]
\centering
\begin{minipage}[t]{.47\linewidth}
\centering
\begin{tikzpicture}[causal,x=1.3cm,y=1.1cm]
\node[ch10node] (ch10f7a-X-0) at (0,2) {$X_{t-3}$};
\node[ch10node] (ch10f7a-X-1) at (1,2) {$X_{t-2}$};
\node[ch10node] (ch10f7a-X-2) at (2,2) {$X_{t-1}$};
\node[ch10node] (ch10f7a-X-3) at (3,2) {$X_{t}$};
\node[ch10node,dashed] (ch10f7a-Z-0) at (0,1) {$Z_{t-3}$};
\node[ch10node,dashed] (ch10f7a-Z-1) at (1,1) {$Z_{t-2}$};
\node[ch10node,dashed] (ch10f7a-Z-2) at (2,1) {$Z_{t-1}$};
\node[ch10node,dashed] (ch10f7a-Z-3) at (3,1) {$Z_{t}$};
\node[ch10node] (ch10f7a-Y-0) at (0,0) {$Y_{t-3}$};
\node[ch10node] (ch10f7a-Y-1) at (1,0) {$Y_{t-2}$};
\node[ch10node] (ch10f7a-Y-2) at (2,0) {$Y_{t-1}$};
\node[ch10node] (ch10f7a-Y-3) at (3,0) {$Y_{t}$};
\draw[ed] (ch10f7a-X-0)--(ch10f7a-X-1);
\draw[ed] (ch10f7a-X-1)--(ch10f7a-X-2);
\draw[ed] (ch10f7a-X-2)--(ch10f7a-X-3);
\draw[ed] (ch10f7a-Y-0)--(ch10f7a-Y-1);
\draw[ed] (ch10f7a-Y-1)--(ch10f7a-Y-2);
\draw[ed] (ch10f7a-Y-2)--(ch10f7a-Y-3);
\draw[ed] (ch10f7a-Z-0)--(ch10f7a-X-1);
\draw[ed] (ch10f7a-Z-1)--(ch10f7a-X-2);
\draw[ed] (ch10f7a-Z-2)--(ch10f7a-X-3);
\draw[ed] (ch10f7a-Z-0)--(ch10f7a-Y-2);
\draw[ed] (ch10f7a-Z-1)--(ch10f7a-Y-3);
\end{tikzpicture}
\par\small (a) 由于存在潜在共同原因$Z$，Granger因果性错误地推断$X$对$Y$存在因果影响。
\end{minipage}\hfill
\begin{minipage}[t]{.47\linewidth}
\centering
\begin{tikzpicture}[causal,x=1.3cm,y=1.15cm]
\node[ch10node] (ch10f7b-X-0) at (0,1) {$X_{t-3}$};
\node[ch10node] (ch10f7b-X-1) at (1,1) {$X_{t-2}$};
\node[ch10node] (ch10f7b-X-2) at (2,1) {$X_{t-1}$};
\node[ch10node] (ch10f7b-X-3) at (3,1) {$X_{t}$};
\node[ch10node] (ch10f7b-Y-0) at (0,0) {$Y_{t-3}$};
\node[ch10node] (ch10f7b-Y-1) at (1,0) {$Y_{t-2}$};
\node[ch10node] (ch10f7b-Y-2) at (2,0) {$Y_{t-1}$};
\node[ch10node] (ch10f7b-Y-3) at (3,0) {$Y_{t}$};
\draw[ed] (ch10f7b-X-0)--(ch10f7b-Y-1);
\draw[ed] (ch10f7b-Y-0)--(ch10f7b-X-1);
\draw[ed] (ch10f7b-X-1)--(ch10f7b-Y-2);
\draw[ed] (ch10f7b-Y-1)--(ch10f7b-X-2);
\draw[ed] (ch10f7b-X-2)--(ch10f7b-Y-3);
\draw[ed] (ch10f7b-Y-2)--(ch10f7b-X-3);
\end{tikzpicture}
\par\small (b) 如果$X_t$对$Y_{t+1}$的影响与$Y_t$对$X_{t+1}$的影响都是确定性的，Granger因果性会错误地推断$X$对$Y$及$Y$对$X$均无因果影响。
\end{minipage}
\caption{在这些例子中，Granger因果性推断出错误的图结构。}
\end{figure}


Ay and Polani［2008］给出了Granger因果性失效的第二个例子，见图10.7(b)。假设$X_{t-1}$通过复制操作确定性地影响$Y_t$，即$Y_t:=X_{t-1}$；同样，将$Y_{t-1}$的值复制给$X_t$。直观上，$X$与$Y$显然强烈相互影响：干预$X_t$的值，会改变所有$Y_{t+1+2k}$、$k\in\mathbb N_0$的值；同样，干预$Y_t$会改变所有$X_{t+1+2k}$的值。然而，在利用$Y$的过去预测$Y_t$时，$X$的过去毫无用处，因为仅凭$Y$自身的过去，就已经能够完美预测$Y_t$。

当然，确定性关系通常会给基于条件独立的因果推断带来困难，因为确定性会诱导额外的独立关系。例如，在因果链$X\to Y\to Z$中，如果$Y$是$X$的函数，就会得到$Y\indep Z\mid X$，而这不是该因果结构的典型性质。因此，人们可能认为这一例子过于人为，更自然的版本应当是带噪声的复制操作。对于$X_t$与$Y_t$均为二值变量的情况，Janzing et al.［2013，例7］证明，当复制操作的噪声水平趋于$0$时，转移熵也趋于$0$。这时，Granger因果性确实会推断$X$导致$Y$、$Y$也导致$X$；但是，当噪声很小时，加入$X$的过去对预测$Y_t$所带来的微小改善，无法恰当地反映两个时间序列之间在直观意义下依然很强的相互影响。

从这一意义上说，转移熵并不是一个时间序列对另一个时间序列的因果影响强度的恰当度量。Janzing et al.［2013］讨论了不同因果影响量化方案的局限，范围包括时间序列与独立同分布情形，并提出另一种基于信息论的因果强度度量。归纳本段讨论，需要强调：对两个满足因果充分性的时间序列，关于因果影响是否存在的\emph{定性}陈述，只会在相当人为的情形中失效；而用转移熵\emph{量化}因果影响——这是将“预测改善”作信息论解释时容易产生的想法——在不那么人为的情形下也可能出问题。

还存在另一种情形：Granger因果性在定量上会产生误导，但只要不违反忠实性，其定性陈述仍然正确。不过，这种情形涉及即时效应，而有人可能认为，只要时间分辨率足够精细，即时效应就会消失［Granger, 1988］。对于图10.8(a)，d-分离给出
\[
 Y_t\indep X_{\mathrm{past}(t)}\mid Y_{\mathrm{past}(t)}.
\]
直观地说，只有当前值$X_t$有助于更好地预测$Y_t$，过去的取值$X_{t-1},X_{t-2},\ldots$则没有用处，因此Granger因果性不会提出从$X$到$Y$的联系。不过，对于图10.8(b)，如果假设忠实性成立，Granger因果性确实能够检测$X$对$Y$的影响，尽管这种影响仍是纯即时的。但是，预测的少量改善，并不能恰当反映$X_t$对$Y_t$可能很强的影响。
\begin{figure}[htbp]
\centering
\begin{minipage}[t]{.47\linewidth}
\centering
\begin{tikzpicture}[causal,x=1.3cm,y=1.15cm]
\node[ch10node] (ch10f8a-X-0) at (0,1) {$X_{t-3}$};
\node[ch10node] (ch10f8a-X-1) at (1,1) {$X_{t-2}$};
\node[ch10node] (ch10f8a-X-2) at (2,1) {$X_{t-1}$};
\node[ch10node] (ch10f8a-X-3) at (3,1) {$X_{t}$};
\node[ch10node] (ch10f8a-Y-0) at (0,0) {$Y_{t-3}$};
\node[ch10node] (ch10f8a-Y-1) at (1,0) {$Y_{t-2}$};
\node[ch10node] (ch10f8a-Y-2) at (2,0) {$Y_{t-1}$};
\node[ch10node] (ch10f8a-Y-3) at (3,0) {$Y_{t}$};
\draw[ed] (ch10f8a-Y-0)--(ch10f8a-Y-1);
\draw[ed] (ch10f8a-Y-1)--(ch10f8a-Y-2);
\draw[ed] (ch10f8a-Y-2)--(ch10f8a-Y-3);
\draw[ed] (ch10f8a-X-0)--(ch10f8a-Y-0);
\draw[ed] (ch10f8a-X-1)--(ch10f8a-Y-1);
\draw[ed] (ch10f8a-X-2)--(ch10f8a-Y-2);
\draw[ed] (ch10f8a-X-3)--(ch10f8a-Y-3);
\end{tikzpicture}
\par\small (a) Granger因果性无法检测$X$对$Y$的影响，因为$X$的过去仅通过$Y$的过去影响$Y_t$。
\end{minipage}\hfill
\begin{minipage}[t]{.47\linewidth}
\centering
\begin{tikzpicture}[causal,x=1.3cm,y=1.15cm]
\node[ch10node] (ch10f8b-X-0) at (0,1) {$X_{t-3}$};
\node[ch10node] (ch10f8b-X-1) at (1,1) {$X_{t-2}$};
\node[ch10node] (ch10f8b-X-2) at (2,1) {$X_{t-1}$};
\node[ch10node] (ch10f8b-X-3) at (3,1) {$X_{t}$};
\node[ch10node] (ch10f8b-Y-0) at (0,0) {$Y_{t-3}$};
\node[ch10node] (ch10f8b-Y-1) at (1,0) {$Y_{t-2}$};
\node[ch10node] (ch10f8b-Y-2) at (2,0) {$Y_{t-1}$};
\node[ch10node] (ch10f8b-Y-3) at (3,0) {$Y_{t}$};
\draw[ed] (ch10f8b-X-0)--(ch10f8b-X-1);
\draw[ed] (ch10f8b-X-1)--(ch10f8b-X-2);
\draw[ed] (ch10f8b-X-2)--(ch10f8b-X-3);
\draw[ed] (ch10f8b-Y-0)--(ch10f8b-Y-1);
\draw[ed] (ch10f8b-Y-1)--(ch10f8b-Y-2);
\draw[ed] (ch10f8b-Y-2)--(ch10f8b-Y-3);
\draw[ed] (ch10f8b-X-0)--(ch10f8b-Y-0);
\draw[ed] (ch10f8b-X-1)--(ch10f8b-Y-1);
\draw[ed] (ch10f8b-X-2)--(ch10f8b-Y-2);
\draw[ed] (ch10f8b-X-3)--(ch10f8b-Y-3);
\end{tikzpicture}
\par\small (b) 这里，$X$的过去仍有助于预测$Y_t$，因为$X_{t-1}$通过$X_t$间接影响$Y_t$。因此，Granger因果性仍能检测$X$对$Y$的影响。
\end{minipage}
\caption{两种存在即时效应的情形：Granger因果性在(a)中无法检测，在(b)中能够检测。}
\end{figure}

\chtennote{本段开头关于“定性陈述仍然正确”的概括，不能涵盖紧接着的图10.8(a)：即使忠实性成立，该图中的即时因果效应也可能无法由Granger因果性检测。后文的正确检测结论针对图10.8(b)。原文两种情形及其结论均予保留。}

为处理即时效应，已有研究提出修改Granger因果性，在对应结构因果模型中加入即时项；但这时可识别性可能失效［例如Lütkepohl, 2007，式(2.3.20)、(2.3.21)］。如果已知系统包含即时效应，一种看似自然的修改是：在式(10.8)中，将$Y_t$不仅对$X_{\mathrm{past}(t)}$回归，而是对$X_t\cup X_{\mathrm{past}(t)}$回归。然而，正如Granger［1988］已经指出的，这可能导致错误结论：$X_t$有助于预测$Y_t$，同样可能意味着$Y_t$影响$X_t$，并不一定表明存在从$X_t$到$Y_t$的影响。

\begin{remark}[模型误设可能有帮助]
这一认识包含一条看似矛盾的信息：即使变量之间存在即时影响，在推断因果陈述时，检验某个变量的\emph{过去}是否有助于预测，也比检验其\emph{过去和当前值}是否有助于预测更有说服力。定理10.4的第(i)项并未排除即时效应。因此，只要满足因果充分性，仍然可以得出结论：在根据$Y_{\mathrm{past}(t)}$预测$Y_t$时，$X_{\mathrm{past}(t)}$带来的任何益处，都源于$X$对$Y$的影响。此外，只要$X$对$Y$存在某种影响，无论是否纯即时，在一般情形下，给定$Y_{\mathrm{past}(t)}$后，$X_{\mathrm{past}(t)}$都会改善对$Y_t$的预测。
\end{remark}
\FloatBarrier
